\documentclass[12pt,a4paper,utf8x]{report}
\usepackage [francais]{babel}
\usepackage[utf8]{inputenc}  
\usepackage[T1]{fontenc} 
% Pour pouvoir utiliser 
\usepackage{ucs}
\usepackage{textcomp}
\usepackage{graphicx}
\usepackage{keystroke}
\usepackage{amssymb}
\usepackage{amsmath}
%\usepackage{pifont}
\usepackage{url} % Pour avoir de belles url
\usepackage{geometry}
\usepackage{hyperref}


\usepackage {listings}% Pour mettre du code source
\lstset{language=sh}

% Pour pouvoir passer en paysage
\usepackage{lscape}

% Pour pouvoir faire plusieurs colonnes
\usepackage {multicol}

\usepackage{makeidx}% Pour crééer un index
\usepackage{graphicx} % Pour insérer des images (entres autres)
\usepackage[cc]{titlepic} %rajouter le logo 	 dans la page de garde
\usepackage{tocbibind}
\usepackage{glossaries} % Créer un glossaire
\hypersetup{
  backref=true,
  %permet d'ajouter des liens dans...
  pagebackref=true,%...les bibliographies
  hyperindex=true, %ajoute des liens dans les index.
  colorlinks=true, %colorise les liens
  breaklinks=true, %permet le retour à la ligne dans les liens trop longs
  urlcolor= blue, %couleur des hyperliens
  bookmarks=true, %créé des signets pour Acrobat
  bookmarksopen=true,
  %si les signets Acrobat sont créés,
  %les afficher complètement.
  pdftitle={PG\_Rapport}, %informations apparaissant dans
  pdfauthor={\textsc{Marguerite} Alain, \textsc{Pluchon} Gwenael\\ \textsc{Rincé} Romain, \textsc{Sebari} Nabila},
  %dans les informations du document
  pdfsubject={PG\_Rapport}
  %sous Acrobat.
}
%Entête pied de page
%Définition des entêtes : 
\usepackage{fancyhdr}
\pagestyle{fancy}

\lstset{language=XML,    numbers=left
   , tabsize=2
   , frame=single
   , breaklines=true
   , basicstyle=\ttfamily
   , numberstyle=\tiny\ttfamily
   , framexleftmargin=13mm
   , xleftmargin=12mm
   %, frameround={tttt}
   , captionpos=b  }
%\usepackage{pifont}
\usepackage{url} % Pour avoir de belles url

% Pour mettre du code source
\usepackage {listings}
% Pour pouvoir passer en paysage
\usepackage{lscape}

% Pour pouvoir faire plusieurs colonnes
\usepackage {multicol}
% POur crééer un index
\usepackage{makeidx}
\makeindex

% Pour l'interligne de 1.5
\usepackage {setspace}
% Pour les marges de la page
%\geometry{a4paper, top=2.5cm, bottom=3.5cm, left=1.5cm, right=1.5cm, marginparwidth=1.2cm}

%\parskip=5pt %% distance entre § (paragraphe)
%\sloppy %% respecter toujours la marge de droite 

% Pour les pénalités :
\interfootnotelinepenalty=150 %note de bas de page
\widowpenalty=1500 %% veuves et orphelines
\clubpenalty=1500

%Pour la longueur de l'indentation des paragraphes
%\setlength{\parindent}{15mm}



%%%% debut macro pour enlever le nom chapitre %%%%
\makeatletter
\def\@makechapterhead#1{%
  \vspace*{50\p@}%
  {\parindent \z@ \raggedright \normalfont
    \interlinepenalty\@M
    \ifnum \c@secnumdepth >\m@ne
        \Huge\bfseries \thechapter\quad
    \fi
    \Huge \bfseries #1\par\nobreak
    \vskip 40\p@
  }}

\def\@makeschapterhead#1{%
  \vspace*{50\p@}%
  {\parindent \z@ \raggedright
    \normalfont
    \interlinepenalty\@M
    \Huge \bfseries  #1\par\nobreak
    \vskip 40\p@
  }}
\makeatother

%%%% fin macro %%%%
\hyphenation{a-pé-rio-dique}
\hyphenation{a-pé-rio-diques}
%Couverture 



\title
{
	\normalsize{ M1 ALMA\\ 
	Université de Nantes\\
	2011-2012}\\
	\vspace{15mm}
	\Huge{Rapport de Projet :\\Programmation générative }
}



\author{\textsc{Marguerite} Alain, \textsc{Pluchon} Gwenael\\ \textsc{Rincé} Romain, \textsc{Sebari} Nabila
	\vspace{45mm}
}
%\titlepic{\includegraphics[scale=1.70]{img/logouniv}     \hspace{2cm} \includegraphics[scale=0.12	]{img/logo}}
\date
{	
	\normalsize{Université de Nantes \\ 2 rue de la Houssinière, BP92208, F-44322 Nantes cedex 03, FRANCE
	\\ 
	\vspace{5mm}	
	}
}

\hyphenation{mé-thodes}

\begin{document}
\renewcommand{\labelitemi}{$\bullet$} 	
\maketitle


\clearpage

\tableofcontents
\clearpage

% Pour avoir un interligne de 1,5
\begin{onehalfspace}
\input{intro}
\input{presentation}
%\chapter{Cas concret : Tas binaire}
\chapter[Concrète]{Choix de la structure concrète%
\chaptermark{Concrète}}



Nous avons recherché la structure concrète la plus performante dans l'ensemble des opérations. Ne pouvant savoir quelle serait la nature de l'utilisation de cette classe, nous avons opté pour une structure offrant les meilleures performances pour l'implémentation du tas.
Plusieurs solutions se présentaient pour une implémentation de tas binaire :
\begin{description}
\item
[Tableau des pères] Non approprié (dédié aux arborescence non-ordonnées).
\item[Une liste d'adjacence] Un noeud étant un élément  (étiquette, père, liste des fils) de la liste. Ses perfomances sont dégradées lors d'opération comme la supression ($O(n)$ pour un tas binaire ) ou l'accès à un noeud à partir de son étiquette ($O(n)$). Dans le cas d'une utilisation où les opérations principales aurait été l'ajout d'un fils par exemple, cette solution aurait pu être retenue.
\item[Liste de fratrie]  Un noeud étant un élément (étiquette, père, liste des frères, premier fils)  de la liste.
\item[Tableau modulaire]
\textbf{definition manquante} Structure présentant les meilleurs complexités pour l'ensemble des opérations. Cette solution répond le mieux aux contraintes que nous avons exposés précedement. C'est donc elle que nous avons retenue. 
\end{description} 


\section[Developpement]{Developpement du tas binaire%
\sectionmark{Developpement}}
\sectionmark{Developpement}

%\section{Developpement du tas binaire}
Une fois la sélection de la structure faite. Nous avons au cours des semaines adopté le plan de développement suivant : 
\begin{enumerate}
\item
Définition des prototypes des méthodes
\item
Implémentation des méthodes
\item
Rajout de l'itérateur 
\item
Rajout de l'allocateur et du comparator
\item
Debug
\end{enumerate}
Cette démarche nous a permis d'avancer rapidément au cours des première semaine. Les cours d'algorithmiques du 1\up{er} semestre encore en tête. Nous avons rapidement déployées les méthodes principales du tas binaire (Extraire, Tasser, Insérer\dots), sans trop nous soucier de la syntaxe C++ (nouvelle pour notre bînome).
Nous avons amèrement regretté notre insoucience de ne jamais tester nos méthodes avant la 5\up{ème} ou 6\up{ème} semaine. Nous avons donc passé un temps important à débuguer un nombre de d'erreurs de syntaxe conséquent.

La création d'un itérateur a été un réel frein dans le développement du projet.\\Notre premier objectif était de construire un itérateur permettant de parcourir les valeurs dans un ordre trié dépendant de la classe \verb+Compare+. Cependant un tel itérateur nécessitait de maintenir une structure de taille égale à celle du tas, ce qui s'avérait coûteux en mémoire.
Nous avons donc décidé de nous orienter vers un itérateur \verb+DFS+ plus classique. Malgré tout nous avons eu du mal a maitriser la syntaxe C++ lié aux itérateurs et ils nous a fallu malgré tout un certain temps.

\chapter[Généralisation]{Généralisation des différentes structures d'arbres%
\chaptermark{Généralisation}}

\input{Conclu}


% Pour finir l'interligne de 1,5
\end{onehalfspace}



\listoffigures

\printindex

\appendix

\bibliographystyle{alpha}
\bibliography{biblio.bib}


\end{document}
